\begin{UseCase}{CU-52}{Registrar contexto indígena o rural del NNA}{
	Permite capturar el nivel de bilingüismo de la familia, la participación en asambleas comunitarias y el parentesco con autoridades locales cuando el NNA vive en comunidad indígena o rural.
}
	\UCitem{Versión}{\color{Gray}1.0}
	\UCitem{Autor}{\color{Gray}Rivera Segura José Emiliano}
	\UCitem{Supervisa}{\color{Gray}Guerra Salinas Edgar Rafael}
	\UCitem{Actor}{\hyperlink{tTSocial}{Trabajador Social}}
	\UCitem{Propósito}{Documentar el contexto cultural e institucional del NNA en comunidades indígenas o rurales para adaptar el PRD a la realidad comunitaria.}
	\UCitem{Entradas}{Nivel de bilingüismo de la familia (texto), participación en asambleas (texto), cargos desempeñados en la comunidad, parentesco con autoridades locales.}
	\UCitem{Origen}{Teclado.}
	\UCitem{Salidas}{Contexto comunitario registrado en el expediente.}
	\UCitem{Destino}{Pestaña Datos Familiares, subsección Contexto comunitario.}
	\UCitem{Precondiciones}{Usuario pertenece al equipo asignado.}
	\UCitem{Postcondiciones}{Datos del contexto comunitario almacenados.}
	\UCitem{Errores}{Ninguno específico.}
	\UCitem{Tipo}{Caso de uso primario}
	\UCitem{Observaciones}{Conforme a §4.3.3 de la Caja de Herramientas.}
\end{UseCase}

\begin{UCtrayectoria}
	\UCpaso[\UCactor] Accede a la subsección \IUbutton{Contexto Comunitario}.
	\UCpaso Despliega el formulario de contexto indígena o rural.
	\UCpaso[\UCactor] Captura el nivel de bilingüismo de la familia, participación comunitaria, cargos y parentesco con autoridades.
	\UCpaso[\UCactor] Presiona \IUbutton{Guardar contexto comunitario}.
	\UCpaso Almacena los datos y muestra \hyperlink{mMSG09}{\bf MSG-09}.
	\UCpaso[] -- -- Fin del caso de uso.
\end{UCtrayectoria}

%-------------------------------------- CU-53
